\section{Method}
\label{sec:method}

\subsection{Distillation of one layer is least squares}
\label{sec:objective}

Let $\theta_{\mathrm{b}}$ be a base model and $\theta_1,\dots,\theta_K$
specialists post-trained from it, one per domain; we call them teachers, and
the model we construct, $\theta_\star$, the student. All share an architecture.
Consider a linear layer with base weight $\Wb\in\mathbb{R}^{m\times n}$ and
teacher weights $W_d=\Wb+\dd_d$. Let $\mu_d$ be the distribution of the input
$x\in\mathbb{R}^n$ this layer receives inside teacher $d$, to be specified in
Section~\ref{sec:stats}. We ask the student's layer to reproduce each teacher's
output on that teacher's inputs,
\begin{equation}
W_\star=\arg\min_{W}\ \sum_{d=1}^{K} w_d\,\mathbb{E}_{x\sim\mu_d}\big\|Wx-W_dx\big\|_2^2 ,
\label{eq:obj}
\end{equation}
with domain weights $w_d>0$. This is the distillation objective restricted to
one layer, with the teacher's activations as inputs. It is quadratic in $W$
and depends on $\mu_d$ only through the second moment
$G_d=\mathbb{E}_{x\sim\mu_d}[xx^{\!\top}]$. Setting the gradient to zero gives
$W_\star\sum_d w_dG_d=\sum_d w_dW_dG_d$, and since the base weight cancels,
\begin{equation}
W_\star=\Wb+\Big(\sum_{d}\dd_d\,\Gt_d\Big)\Big(\sum_{d}\Gt_d\Big)^{-1},
\label{eq:solve}
\end{equation}
where $\Gt_d=w_dG_d$. This is the RegMean solution \citep{jin2023regmean}
written for updates. Following that work we shrink off-diagonal second moments,
\begin{equation}
\Gt_d=w_d\big[\alpha\,G_d+(1-\alpha)\,\mathrm{diag}(G_d)\big],\qquad \alpha=0.9,
\label{eq:shrink}
\end{equation}
which makes the system well conditioned and is equivalent to adding a
per-coordinate penalty that pulls $W$ toward each teacher
(Appendix~\ref{app:derivation}).

Equation~\eqref{eq:solve} makes the contrast with weight arithmetic explicit.
Task arithmetic sets the update to $\sum_d\dd_d$ and so applies every
teacher's update to every input, including inputs that only another teacher
ever sees; on teacher $e$'s inputs the layer output is off by
$\sum_{d\ne e}\dd_dx$. \method minimizes exactly that error, weighted by how
often each input direction occurs. Three properties of the solution locate
where teachers can conflict.

\begin{proposition}
\label{prop:conflict}
Let $\dd_\star=W_\star-\Wb$ solve Equation~\eqref{eq:solve} with
$\Gt_d=w_dG_d$.
\begin{enumerate}
\item[(i)] If $G_d=G$ is invertible for all $d$, then
$\dd_\star=\sum_dw_d\dd_d/\sum_dw_d$.
\item[(ii)] If $G_e$ is invertible and $w_d/w_e\to0$ for all $d\ne e$, then
$\dd_\star\to\dd_e$.
\item[(iii)] If the ranges of $G_1,\dots,G_K$ are mutually orthogonal and the
inverse is taken as the pseudo-inverse, then $\dd_\star x=\dd_ex$ for every
$x$ in the range of $G_e$ and every $e$.
\end{enumerate}
\end{proposition}

Appendix~\ref{app:derivation} gives the proofs. By (i), the closed form reduces
to weight averaging, not to task arithmetic, when the teachers cannot be told
apart by their inputs. By (iii), interference arises only where input
distributions overlap; Equation~\eqref{eq:diag} below measures how much.

\subsection{On-policy statistics}
\label{sec:stats}

The solution is only as good as $\mu_d$. We take $\mu_d$ to be the distribution
of layer inputs when teacher $d$ processes its own answers to problems from
its domain. This is the distribution the layer operates on in use, since a
reasoning model's context is dominated by its own generation, and it cannot be
obtained from reference solutions, which the model did not write. We refer to
statistics gathered on model-generated rollouts as on-policy, in contrast to
statistics gathered on a fixed corpus; the rollouts in our main configuration
are the teachers', as in sequence-level distillation \citep{kim2016sequence},
and Section~\ref{sec:ablations} evaluates the variant that uses the student's.

Concretely, teacher $d$ answers a set of training problems from its domain,
sampled exactly as in evaluation. After selection
(Section~\ref{sec:selection}), the teacher reads each kept sequence, prompt
and answer, in a single teacher-forced forward pass. A forward pre-hook on
every linear layer of every decoder block accumulates $X^{\!\top}X$ over all
token positions, and $G_d$ is that sum divided by the number of tokens, so
that domains with longer answers do not dominate by length. Layers that read
the same tensor share a Gram matrix: in the architectures we use, the query,
key and value projections form one group, the gate and up projections another,
and the attention output and down projections one each. The grouping is
detected automatically from one short forward pass, and the same statistics
are collected for the final hidden states to solve the output head.

\subsection{Solving}
\label{sec:solve}

For each group, Equation~\eqref{eq:solve} is solved once per member layer in
double precision. Tensors that are not linear maps of a hidden state ---
embeddings, normalization gains and biases --- change little and are combined
by task arithmetic. Besides the fused model, the solve returns for every layer
and teacher the share of that teacher's update left unreproduced,
\begin{equation}
\varepsilon_d=\frac{\tr\!\big[(\dd_\star-\dd_d)\,\Gt_d\,(\dd_\star-\dd_d)^{\!\top}\big]}
{\tr\!\big[\dd_d\,\Gt_d\,\dd_d^{\!\top}\big]},
\label{eq:diag}
\end{equation}
with $\Gt_d$ of Equation~\eqref{eq:shrink}. It is zero when the fused layer
matches the teacher on its inputs, and it locates the layers where the
teachers compete. Figure~\ref{fig:overview} illustrates the procedure, and
Algorithm~\ref{alg:main} in Appendix~\ref{app:eval} states it in full.

\subsection{Which rollouts to use}
\label{sec:selection}

Every token of every kept answer contributes one outer product to $G_d$, so
the statistics are a token-weighted picture of the teacher's behavior. That
picture is distorted by a failure mode of long-form generation. Some answers
do not terminate: the model enters a repetition loop and continues to the
length limit. Such answers are few, but each is as long as the limit allows,
and the layer inputs inside a loop are nearly periodic, so their contribution
to the Gram matrix is large and concentrated on a few directions. The
least-squares solution then spends its capacity reproducing the teacher on
degenerate input.

We detect loops without supervision from the compressibility of the text: an
answer is flagged when the last 20{,}000 characters compress, under DEFLATE, to
less than 5\% of their length, far below ordinary reasoning text. Our default
selection keeps answers that finished on their own and are not flagged. Where
answers can be graded automatically, we additionally keep only correct ones:
mathematics answers are checked against the reference answer, and code
answers are run against the tests of their problem.

\subsection{Precision}
\label{sec:precision}

Released checkpoints are stored in bfloat16, which keeps eight significant
bits. Reinforcement learning with small learning rates barely moves weights at
that resolution, by about one representable step
(Section~\ref{sec:exp-precision}). A fused update is no larger, and unlike a
specialist's it is not made of whole steps. We therefore store fused models
in 32-bit floats, since rounding the fused update to bfloat16 measurably
distorts it. A fused model is evaluated in float16, which keeps three more
significant bits than bfloat16 and fits the memory of a single GPU where
float32 does not. A format change by itself shifts measured accuracy, so we
state the format behind every score.
